\subsection{A uniform evaluated tail for every inert first prime}

\begin{statement}{E.114}{Proved}{All S-class and cyclic unit norms in the inert-prime positional tail}
Let $\ell\ne3$ be a rational prime with $\ell\equiv2\pmod3$,
and put
\[
 \begin{gathered}
 s=v_3(\ell+1),\qquad t=3^{s-1},\qquad
 K_j=\mathbb Q(\xi_j),\quad \xi_j=\zeta_{3^j},\\
 L^{\ell}_{m,j}=K_j(\beta_{m,j}),\qquad
 \beta_{m,j}=\ell^{1/3^m}\xi_{m+j},\quad j\ge s, m\ge0.
 \end{gathered}
 \tag{E114.1}
\]
The roots of $\ell$ are positive and compatible, and the
roots of unity are those already fixed. Let S consist of
the primes over three and $\ell$. There is one three-prime
and $t$ primes over $\ell$ in every field in (E114.1).
All are totally ramified in each vertical cubic step.
The entire three-primary S-class system is
\[
           \operatorname{Cl}_S(L^{\ell}_{m,j})\{3\}=0
                    \quad(j\ge s, m\ge0),
 \tag{E114.2}
\]
with the actual ideal norm and extension maps between
these zero groups.

Put $U_{m,j}=\mathcal O_{L^{\ell}_{m,j},S}^{\times}$ and
$r_{m,j}=3^m\varphi(3^j)/2$. For every $1\le r\le j$,
the actual extension $L^{\ell}_{m+r,j}/L^{\ell}_{m,j}$
is cyclic of degree $n=3^r$. Its complete unit-norm
quotient and completed norm kernel are
\[
 \begin{gathered}
 U_{m,j}/N U_{m+r,j}\xrightarrow{\sim}(\mathbb Z/n)^t,
 \qquad u\longmapsto
   \bigl(\log_{\zeta_n}(\beta_{m,j},u)_{v,n}\bigr)_{v\mid\ell},\\
 \ker(N:\widehat U_{m+r,j,3}\longrightarrow
                          \widehat U_{m,j,3})
       \simeq C_n\oplus\mathbb Z_3^{(3^r-1)r_{m,j}}.
 \end{gathered}
 \tag{E114.3}
\]
The finite kernel column is the literal $\mu_n$.
The omitted three-place character is the negative sum
of the displayed $t$ characters. The same quotient map
is exact after completion. Equations (E114.4)--(E114.6)
and (E114.15) give its literal finite-field evaluation and inverse
coordinates, so its zero condition is both necessary
and sufficient for an S-unit norm.
Norms at larger distances are the compositions of the
proved adjacent maps; the cyclic assertion in (E114.3)
has exactly the stated range $r\le j$.

The uniform cubic inverse has a small algebraic form.
Let $\Gamma=\operatorname{Gal}(K_s/K_1)=\langle\sigma\rangle$,
where $\sigma(\xi_s)=\xi_s^4$; its order is $t$.
Choose a prime $\mathfrak t_0\mid\ell$ of $K_s$ and a
completed principal generator $\pi_0$ of that prime,
and put $\pi_g=\sigma^g\pi_0$ and
$\mathfrak t_g=\sigma^g\mathfrak t_0$, $g\in\mathbb Z/t$.
The proof constructs these completed generators from
actual principal ideal powers of order prime to three.
Let
\[
 R_t=\mathbb F_3[\Gamma],\quad I=(\sigma-1)R_t,
 \qquad c=\frac{\ell^2-1}{3^s}\pmod3\in\mathbb F_3^\times,
 \qquad
 a=\sum_{h\in\mathbb Z/t}
       \log_\zeta(\ell,\pi_0)_{\mathfrak t_h,3}\,\sigma^h.
 \tag{E114.4}
\]
Then $a\in I$, and the actual cubic character matrix
on the included $\pi_g$ columns is multiplication by
$c+a$ for $j=s$, and by $c$ for $j>s$. Its full inverse is
\[
            (c+a)^{-1}
             =c^{-1}\sum_{d=0}^{t-1}(-c^{-1}a)^d
                      \quad\text{in }R_t.
 \tag{E114.5}
\]
The identification of a coefficient vector with a
product of the corresponding $\pi_g$ specifies the
section of the quotient (E114.3) at $r=1$.
No value of an uncomputed class group occurs in this
inverse.

For the higher characters in (E114.3) choose an actual
integer $d$ prime to three and $x_0\in K_s^\times$
with $(x_0)=\mathfrak t_0^d$, and take
$\pi_0=x_0^{1/d}$ in the multiplication completion.
Set $x_g=\sigma^g x_0$. At a labeled prime $v_h\mid\ell$
in $K_j$, let $\xi_{j,h}$ be the residue of $\xi_j$,
and put $q_j=\ell^{2\cdot3^{j-s}}$.
The inverse matrix is obtained from the completely
specified $t$ by $t$ matrix
\[
 \begin{split}
 C_{j,r}(h,g)
 &=d^{-1}\log_{\xi_{j,h}^{3^{j-r}}}
    \left(
      \xi_{j,h}^{d\delta_{hg}}
       \overline{x_g/\ell^{d\delta_{hg}}}^{\,-1}
    \right)^{(q_j-1)/3^r}
                              \quad\text{in }\mathbb Z/3^r,\\
 (z_h)_h&\longmapsto
     \prod_g\pi_g^{\,(C_{j,r}^{-1}z)_g}
                  \quad\text{in }\widehat U_{m,j,3}/N\widehat U_{m+r,j,3}.
 \end{split}
 \tag{E114.6}
\]
Its reduction modulo three is the invertible matrix
just proved. Thus its inverse is computed over
$\mathbb Z/3^r$ by ordinary finite matrix inversion,
or by lifting (E114.5) through
$B\mapsto B(2-C_{j,r}B)$.
For actual, uncompleted unit representatives replace
$1/d$ in the exponents by its inverse modulo $3^r$.

Let
$P_j(\ell)=\{p\ne3,\ell:v_3(\operatorname{ord}_{\mathbb F_p^\times}\ell)=j\}$.
The entire completed norm factor in each tail is
\[
 \begin{gathered}
 \mathcal X_{\ell,j}=\varprojlim_{m,N}
                       \widehat{(L^{\ell}_{m,j})^\times}_3,
 \qquad
 \mathcal U_{\ell,j}=\varprojlim_{m,N}\widehat U_{m,j,3},\\
 0\longrightarrow\mathcal U_{\ell,j}
   \longrightarrow\mathcal X_{\ell,j}
   \xrightarrow{\operatorname{div}_{\notin S}}
          \mathcal R(P_j(\ell))\longrightarrow0,
 \qquad\mathcal U_{\ell,j}\simeq\prod_{i\ge1}\mathbb Z_3.
 \end{gathered}
 \tag{E114.7}
\]
The sequence splits by choosing coherent completed
S-principal generators $a_{j,p}$ of the unique surviving
degree-one prime chains and putting
$s_j(v)=\prod_pa_{j,p}^{v_p}$.
Its inverse is $x\mapsto(x/s_j(\operatorname{div}x),\operatorname{div}x)$.
The compact coordinate inverse is the character-basis
construction of (E110.7), using the actual S-unit norms
in (E114.3). All products and quotient topologies are
those of multiplication completion, with restricted
prime exponents as in (E110.2).

Here is the exact rational norm retained by these maps.
For $y\in\mathcal U_{\ell,j}$ read its initial $K_j$
coordinate and put
\[
 \begin{gathered}
 \nu_3(y)=v_{\mathfrak q_3}(y_0),\qquad
 \nu_\ell(y)=2\cdot3^{j-s}\sum_{v\mid\ell}v(y_0),\\
 N(y s_j(v))=
 3^{\nu_3(y)+\sum\rho_{j,p}v_p}
 \ell^{\nu_\ell(y)+\sum\eta_{j,p}v_p}
                       \prod_{p\in P_j(\ell)}p^{v_p}.
 \end{gathered}
 \tag{E114.8}
\]
Here $\rho_{j,p},\eta_{j,p}$ are the same initial
valuation maps applied to the chosen $a_{j,p}$.
They belong to $\mathbb Z_3$, so all sums converge.

The comparison to global sections has a literal split
map even when the lower positional factors have not
been evaluated. Fix $k\ge s-1$ and the first cocycle
$A=3^{k+1}\kappa_3(\ell)$.
At its cofinal positional levels the factors with
$j\ge s$ are exactly (E114.1). Let $\mathcal N_{<s}$
be the actual norm system formed from its pure factor
and its factors with $1\le j<s$.
Then the geometric-conjugacy classes of sections in
$H$ have the following proved topological product:
\[
 \begin{gathered}
 \mathscr S_{3^k\kappa(\ell)}(H)
   \simeq
   \left(\prod_{j=s}^{k+1}\mathcal R(P_j(\ell))\right)
                                      \times\mathcal K_{\ell,k},\\
 \mathcal K_{\ell,k}=
 \left\{(z,(y_j))\in\mathcal N_{<s}\times
                \prod_{j=s}^{k+1}\mathcal U_{\ell,j}:
     N(z)\prod_{j=s}^{k+1}3^{\nu_3(y_j)}\ell^{\nu_\ell(y_j)}
                     \in(\widehat{\mathbb Q^\times}_3)^3\right\}.
 \end{gathered}
 \tag{E114.9}
\]
The forward map reads $v_j/3$ from the tail divisors
and removes the actual norm-coherent tuple
$s_j(v_j)$. The inverse multiplies back those tuples.
Thus the first product is an explicitly split closed
subgroup of the complete section fiber, obtained by
taking all lower components and all $y_j$ equal to one
and using $s_j(3v_j)$ in each tail.
Its scaled second toric coordinate is
\[
   \left(\sum_{j,p}\rho_{j,p}v_{j,p}\right)\kappa_3(3)
  +\left(\sum_{j,p}\eta_{j,p}v_{j,p}\right)\kappa_3(\ell)
  +\sum_{j,p}v_{j,p}\kappa_3(p).
 \tag{E114.10}
\]
These are actual global $H$ sections, through E.46's
Kummer--Shapiro map. Formula (E114.9) retains the exact
remaining lower norm system and its norm coupling;
it does not assign it the zero S-class calculation
proved only for $j\ge s$. The maps from intrinsic
metabelian or marked sections to these classes remain
the actual quotient maps of E.38.
Every set $P_j(\ell)$ occurring here is infinite.
For a three-adic unit $u$, the arithmetic automorphism
$\Phi_u(a,v(T))=(ua,v(T^u))$ transports (E114.9)
to the first class $3^ku\kappa_3(\ell)$, preserving
the second toric coordinate and with inverse $\Phi_{u^{-1}}$.
The manuscript anchor is \S51.16, scans 172--178,
folios 733v--736v, equations (371)--(379), and Part E(c),(d).
\end{statement}

\begin{proof}
\emph{Fields, local degrees and principal prime columns.}
For $j\ge s$, the binomial expansion gives
\[
 \operatorname{ord}_{3^j}(\ell)=2\cdot3^{j-s},
 \qquad v_3(\ell^{2\cdot3^a}-1)=s+a.
 \tag{E114.11}
\]
There are consequently $t$ primes over $\ell$ in $K_j$,
with the residue degrees in (E114.8), and one prime over
three. At $\ell$ the radicand $\ell\xi_j$ has valuation
one, so every vertical field is totally ramified there.
At three, shifting $X^{3^m}-\ell\xi_j$ by $X=Y-1$
gives constant term $-(\ell+1)-\ell(\xi_j-1)$ of
$1-\xi_j$ valuation one. All its other nonleading
coefficients are divisible by three. Thus it is
Eisenstein there as well. The polynomial discriminant
shows that no other finite prime ramifies.
The equality $\xi_j=\beta_{m,j}^{3^m}/\ell$ gives the
same explicit rational cyclotomic factor as (E111.13),
with seventeen replaced by $\ell$.

The proved one-prime cyclotomic induction in E.47 or
E.111 gives $\operatorname{Cl}(K_j)\{3\}=0$.
In particular the ordinary class of $\mathfrak t_0$
has an order $d$ prime to three. Choose an actual
generator $x_0$ of $\mathfrak t_0^d$ and use the
unique multiplication-by-$d$ inverse in the completed
group to define $\pi_0$. Its divisor is exactly
$\mathfrak t_0$. The Galois orbit of these elements
has the valuation basis prescribed in (E114.4).
When an actual unit is needed in a finite quotient,
$x_g^e$ with $ed\equiv1\pmod{3^r}$ is a representative.
This is the exact passage between actual and completed
prime generators, including their norm characters.

\emph{The uniform cubic character inverse.}
The group $\Gamma$ acts regularly on the $t$ primes
over $\ell$, because $\ell\equiv-1\pmod{3^s}$ and
the decomposition group in $\operatorname{Gal}(K_s/\mathbb Q)$
is its order-two subgroup. All elements of $\Gamma$
fix $\mu_3$. Naturality of the Hilbert pairing gives
\[
 \log_\zeta(\ell,\pi_g)_{\mathfrak t_h,3}
       =\log_\zeta(\ell,\pi_0)_{\mathfrak t_{h-g},3}.
 \tag{E114.12}
\]
Thus the actual matrix is multiplication by $a$ in
the group-ring basis, with exactly the displayed
orientation of the indices.

For $s\ge2$, the rational number $\ell$ is a cube in
$\mathbb Q_3$. One explicit Hensel construction starts
with $x=-1+3b$, where
$b\equiv(\ell+1)/9\pmod3$. Then $x^3\equiv\ell\pmod{27}$
and $v_3(3x^2)=1$, so Hensel's inequality is strict.
Therefore the three-place symbol $(\ell,u)_3$ is one
for every local $u$. All symbols of an S-unit at primes
outside S are also one. Global reciprocity implies
$\sum_h\log_\zeta(\ell,\pi_0)_{\mathfrak t_h,3}=0$.
For $s=1$ there is one prime and we take $\pi_0=\ell$.
The cubic pairing is alternating, so $(\ell,\ell)_3=1$
at that prime as well. These arguments prove $a\in I$
in every case in (E114.4).

The other part of the radicand, $\xi_s$, is a residue
unit. At its $\ell$-adic places, arithmetic tame
reciprocity on the prime generator is
\[
             (\xi_s,\pi_g)_{\mathfrak t_h,3}
                      =\zeta^{c\delta_{hg}}.
 \tag{E114.13}
\]
Indeed its residue field has order $\ell^2$, and the
positive valuation convention gives
$\xi_s^{(\ell^2-1)/3}=\zeta^c$.
The same exponent holds at every labeled place, since
$\Gamma$ fixes $\zeta$. Bilinearity now makes the
matrix for $\ell\xi_s$ equal to $a+c$.

In characteristic three,
$R_t=\mathbb F_3[\sigma-1]/((\sigma-1)^t)$.
Since $a\in I$, its $t$th power is zero. Multiplying
the finite geometric series in (E114.5) by $c+a$
leaves $1-(-c^{-1}a)^t=1$. This proves both inverse
maps, not only a determinant assertion. The last
character, at three, is the negative sum of the tame
coordinates. Hence the full sum-zero character
hyperplane has dimension $t$ and is attained by
these explicit columns.

For example, at $\ell=17$ the actual $\sigma$-orbit
of E.111's first prime generator is
$\pi_1,\xi_2^6\pi_3,\pi_2$.
In this order the retained polynomial/unit calculation
gives $a=2(1+\sigma+\sigma^2)$ and $c=2$.
Thus the cubic matrix has diagonal entries one and
off-diagonal entries two; its inverse has diagonal
entries zero and off-diagonal entries one. At $j>2$
the matrix is $2I$. The certificate also evaluates
the literal ninth and twenty-seventh characters,
their inverse matrices, and the common-root norm
comparison, retaining the actual Galois orbit map.

For $j>s$ the actual norm at each labeled completion is
\[
 N_{K_j/K_s}(\ell\xi_j)=\ell^{3^{j-s}}\xi_s,
 \qquad N_{L^{\ell}_{m,j}/K_j}(\beta_{m,j})=\ell\xi_j.
 \tag{E114.14}
\]
There is exactly one prime above each predecessor, so
these equalities are the local norms too. The common-root
cup-product projection formula and
$\operatorname{inv}\operatorname{Cor}=\operatorname{inv}$
preserve each character of an included prime column,
as in (E111.20). At cubic coefficients the first
radicand is $\xi_s$ modulo cubes when $j>s$.
Consequently its matrix is exactly $c$ rather than
$a+c$. The second equality preserves these matrices
through every vertical height. These comparisons use
the proved identities of Milne III, Sections 2--4
\cite{E_MilneCFT} and retain the common root in each
coefficient group.

\emph{All ideal and S-unit norms.}
We apply the actual sequences (E111.21)--(E111.25),
now with $t+1$ totally ramified finite S-places.
Their divisor maps and inverses do not change:
outside S, invariant divisors are extended divisors,
and the kernel of ideal norm is the image of
$\sigma-1$. Hilbert 90 identifies S-capitulation with
$H^1(G,U_E)$ and identifies the invariant S-class
quotient by extended classes with
$(U_F\cap NE^\times)/NU_E$.

At the initial field $K_j$ the three-class group is
zero. At an adjacent induction step this makes the
S-capitulation kernel, which is killed by three,
zero. The S-unit Herbrand quotient is $3^t$:
the $t+1$ fixed finite places give that many trivial
summands, the product formula removes one, and the
complex places give regular summands. Milne VII
Proposition 3.1 and its complete lattice proof give
the same value directly \cite{E_MilneCFT}.
Thus $[U_F:NU_E]=3^t$.

The actual local characters just proved already have
image of that order. The cyclic Hasse norm theorem
identifies their kernel with $U_F\cap NE^\times$,
so this group equals $NU_E$. The exact S-class map
then gives zero invariant upper three-class group.
A nonzero finite three-group module for a three-group
has nonzero fixed points, so the entire upper
three-class group vanishes. Induction proves (E114.2)
and the cubic case of (E114.3), with the actual norm
zero criterion. The source statements used here are
the opened Hasse norm theorem and the explicit
ideal/Hilbert 90 proof of the ambiguous sequence
\cite[VIII, Theorem 3.1]{E_MilneCFT}
\cite[Theorem 1 and Proposition 1]{E_LemmermeyerAmbiguous}.

For $r\le j$, $\mu_{3^r}$ belongs to the lower field,
so the extension of degree $n=3^r$ is actually cyclic.
The $\ell$-place character on an arbitrary S-unit is
computed by using $\beta_{m,j}$ as uniformizer:
\[
 \log_{\zeta_n}(\beta_{m,j},u)_{v,n}
       =-\log_{\overline{\zeta_n}}
          \left(\overline{u/\beta_{m,j}^{v(u)}}
                                             \right)^{(q_j-1)/n}.
 \tag{E114.15}
\]
The tame sign disappears because $n$ is odd.
For the included actual generator $x_g$ its valuation
is $3^m d\delta_{hg}$, and
$\beta_{m,j}^{3^m}=\ell\xi_j$.
Substitution gives exactly (E114.6), independently
of $m$. Reduction of these characters modulo three
is the already proved cubic matrix. Its inverse
therefore lifts over $\mathbb Z/3^r$, as stated.
The characters have image $(\mathbb Z/n)^t$.

The same ideal Hilbert 90 sequence has zero S-class
defect, since both class groups are now evaluated.
Alternatively the degree-$n$ Herbrand quotient is
$n^t$ and the capitulation group is zero. These actual
maps show that the full character kernel is precisely
the S-unit norm image. Completion preserves this
finite quotient of finitely generated groups.
The free ranks of the source and target S-unit groups
are $3^r r_{m,j}+t$ and $r_{m,j}+t$.
There are exactly $\mu_{3^j}$ as three-power roots
in either field: a larger cyclotomic root would give
a nontrivial unramified extension at $\ell$, which
cannot lie in the totally ramified vertical field.
Norm raises these roots to the $n$th power. This proves
the exact torsion and free-rank kernel in (E114.3).

\emph{The actual norm factors and their topology.}
At $p\ne3,\ell$, write the residue of $\ell$ as its
three-primary part times its prime-to-three part.
Its three-primary part has order $3^j$ precisely
for $p\in P_j(\ell)$; the unique surviving embedding
has $\xi_j$ equal to its inverse. The prime-to-three
part then has unique compatible cube roots. These
give the unique degree-one chain in (E114.7).

To verify that no other divisor coordinate survives,
use the actual affine Frobenius root action of E.46.
Its sole possible finite orbit is its fixed point.
The nonfixed pro-three orbit neighborhoods subdivide
into arbitrarily many equal translates, so invariant
$\mathbb Z_3$ measures there vanish. The map from
prime divisors to orbit sums intertwines positional
pushforward with ideal norm, since the ratio of orbit
lengths is the residue degree. This is the exact
comparison proved in E.111, with $\ell$ in place of
seventeen. It gives the stated restricted support
and its inverse fixed-prime map.

Since the finite S-class groups have order prime
to three, every completed outside-S divisor is
principal. Its fiber is a compact completed S-unit
torsor. Compatible infinite divisors have preimages
by the finite intersection argument: solve any
finite set of equations by choosing a principal
element at its highest level and taking norms.
Compactness of the fixed fibers supplies a coherent
preimage. This proves both directions of (E114.7).
The product splitting and its exact change when
$a_{j,p}$ is multiplied by a coherent S-unit are
identical to (E111.7), and restricted exponents
give the required uniform convergence.

The compact kernel is metrizable and torsion free:
all of its torsion coordinates would lie in
$\mu_{3^j}$, and repeated norm cubing kills every
fixed coordinate. Its character group is therefore
countable and divisible. Choosing a basis of the
characters killed by three and successive divisions
gives the precise product map (E110.7); the finite
character equations give its unique continuous
inverse. The finite representatives are S-unit
characters under actual transposed norm maps.
There are infinitely many coordinates: each
$L^{\ell}_{m,j}$ has one three-prime, so Kato's
Theorem 1.1(2) gives a finite-index projection of the
full false-Tate coherent three-unit system into
its three-unit group of rank $r_{m,j}$.
Restriction to the literal chain (E114.1) is an
actual norm-system map, and these ranks are
unbounded \cite{E_KatoUniversalNorms}.
Finally initial ideal norm weights are one at
three and $2\cdot3^{j-s}$ at the $\ell$-primes.
They prove (E114.8), with convergent restricted
corrections and no uncomputed scalar replacement.

\emph{Comparison with the complete fixed-first section set.}
The positional factorization for $3^k\kappa_3(\ell)$
is exactly (E111.13), with seventeen replaced by
$\ell$. Separate its factors into $j<s$ and $j\ge s$.
The surviving rational prime sets are disjoint, so
the total-norm cube condition forces every tail
divisor to be $3v_j$ for a unique
$v_j\in\mathcal R(P_j(\ell))$.
Conversely $s_j(3v_j)=s_j(v_j)^3$, so the tuple
with these tail components and all other components
one has cube total norm and is an actual section
class by E.46. Its divided divisor projection is
exactly $(v_j)$.

Removing this tuple from any section leaves the
tail units $y_j$ and the lower tuple $z$.
Their total rational norm condition is precisely
the one defining $\mathcal K_{\ell,k}$ in (E114.9),
by (E114.8). Multiplication supplies the inverse.
These maps are continuous and mutually inverse.
The split subgroup is closed: the composite of
its inclusion and the continuous divided-divisor
projection is a continuous idempotent on the
Hausdorff section group, and its fixed set is
exactly that subgroup. Dividing its rational norm
by three gives (E114.10), retaining both S-prime
corrections. This establishes the exact global
comparison and leaves only the explicitly displayed
lower norm coupling in that comparison.

Finally, the infinite supply of independent tail
directions follows by an actual Chebotarev class.
In $\mathbb Q(\mu_{3^{j+1}},\ell^{1/3^j})$ the
Kummer extension over the cyclotomic base has degree
$3^j$, as seen from the valuation one at $\ell$.
The automorphism that sends $\xi_{j+1}$ to
$\xi_{j+1}^{1+3^j}$ and sends $\ell^{1/3^j}$ to
$\xi_j\ell^{1/3^j}$ is therefore defined. Its
Frobenius conjugacy class forces
$v_3(p-1)=j$ and makes $\ell$ a noncube modulo $p$;
hence $v_3(\operatorname{ord}_p\ell)=j$.
Chebotarev's theorem, with its complete proof in
Milne VIII, Theorem 7.4, gives infinitely many such
primes \cite{E_MilneCFT}. Their divided divisor
coordinates separate the section classes in the
split subgroup. Substituting $T\mapsto T^u$ in
the actual law (E46.1) proves the stated unit-part
transport and its inverse.
\end{proof}
